\documentclass{article}
\usepackage{physics}
\usepackage{graphicx}
\usepackage{caption}
\usepackage{amsmath}
\usepackage{bm}
\usepackage{framed}
\usepackage{authblk}
\usepackage{empheq}
\usepackage{amsfonts}
\usepackage{esint}
\usepackage[makeroom]{cancel}
\usepackage{dsfont}
\usepackage{centernot}
\usepackage{mathtools}
\usepackage{subcaption}
\usepackage{bigints}
\usepackage{amsthm}
\theoremstyle{definition}
\newtheorem{lemma}{Lemma}
\newtheorem{defn}{Definition}[section]
\newtheorem{prop}{Proposition}[section]
\newtheorem{rmk}{Remark}[section]
\newtheorem{thm}{Theorem}[section]
\newtheorem{exmp}{Example}[section]
\newtheorem{prob}{Problem}[section]
\newtheorem{sln}{Solution}[section]
\newtheorem*{prob*}{Problem}
\newtheorem{exer}{Exercise}[section]
\newtheorem*{exer*}{Exercise}
\newtheorem*{sln*}{Solution}
\usepackage{empheq}
\usepackage{tensor}
\usepackage{xcolor}
%\definecolor{colby}{rgb}{0.0, 0.0, 0.5}
\definecolor{MIT}{RGB}{163, 31, 52}
\usepackage[pdftex]{hyperref}
%\hypersetup{colorlinks,urlcolor=colby}
\hypersetup{colorlinks,linkcolor={MIT},citecolor={MIT},urlcolor={MIT}}  
\usepackage[left=1in,right=1in,top=1in,bottom=1in]{geometry}

\usepackage{newpxtext,newpxmath}
\newcommand*\widefbox[1]{\fbox{\hspace{2em}#1\hspace{2em}}}

\newcommand{\p}{\partial}
\newcommand{\R}{\mathbb{R}}
\newcommand{\C}{\mathbb{C}}
\newcommand{\lag}{\mathcal{L}}
\newcommand{\nn}{\nonumber}
\newcommand{\ham}{\mathcal{H}}
\newcommand{\M}{\mathcal{M}}
\newcommand{\I}{\mathcal{I}}
\newcommand{\K}{\mathcal{K}}
\newcommand{\F}{\mathcal{F}}
\newcommand{\w}{\omega}
\newcommand{\lam}{\lambda}
\newcommand{\al}{\alpha}
\newcommand{\be}{\beta}
\newcommand{\x}{\xi}

\newcommand{\G}{\mathcal{G}}

\newcommand{\f}[2]{\frac{#1}{#2}}

\newcommand{\ift}{\infty}

\newcommand{\lp}{\left(}
\newcommand{\rp}{\right)}

\newcommand{\lb}{\left[}
\newcommand{\rb}{\right]}

\newcommand{\lc}{\left\{}
\newcommand{\rc}{\right\}}


\newcommand{\V}{\mathbf{V}}
\newcommand{\U}{\mathcal{U}}
\newcommand{\Id}{\mathcal{I}}
\newcommand{\D}{\mathcal{D}}
\newcommand{\Z}{\mathcal{Z}}

%\setcounter{chapter}{-1}


\usepackage{enumitem}



\usepackage{listings}
\captionsetup[lstlisting]{margin=0cm,format=hang,font=small,format=plain,labelfont={bf,up},textfont={it}}
\renewcommand*{\lstlistingname}{Code \textcolor{violet}{\textsl{Mathematica}}}
\definecolor{gris245}{RGB}{245,245,245}
\definecolor{olive}{RGB}{50,140,50}
\definecolor{brun}{RGB}{175,100,80}

%\hypersetup{colorlinks,urlcolor=colby}
\lstset{
	tabsize=4,
	frame=single,
	language=mathematica,
	basicstyle=\scriptsize\ttfamily,
	keywordstyle=\color{black},
	backgroundcolor=\color{gris245},
	commentstyle=\color{gray},
	showstringspaces=false,
	emph={
		r1,
		r2,
		epsilon,epsilon_,
		Newton,Newton_
	},emphstyle={\color{olive}},
	emph={[2]
		L,
		CouleurCourbe,
		PotentielEffectif,
		IdCourbe,
		Courbe
	},emphstyle={[2]\color{blue}},
	emph={[3]r,r_,n,n_},emphstyle={[3]\color{magenta}}
}

\newcommand{\diag}{\text{diag}}
\newcommand{\psirot}{\ket{\psi_\text{rot}(t)} }
\newcommand{\RWA}{\ham_\text{rot}^\text{RWA}}


\begin{document}
\begin{framed}
\noindent TA: \textbf{Huan Q. Bui}\\
Course: \textbf{8.421 - AMO I}\\
Problem set: \textbf{\#2 -- UPDATED SOLUTION}\\
Due: Feb 21, 2025.
\end{framed}
	
	

\noindent \textbf{1. Atomic Units.}\\

\noindent Let the atomic unit of $E$ field be $E_A = e/a_0^2$, the field of the ground state of the electron at the site of the proton in hydrogen. In the following we use cgs units. To convert to SI, replace $e^2 \to e^2/4\pi\epsilon_0$. Factors of 2 are ignored.
\begin{enumerate}[label=\alph*)]
	\item An electron of mass m is confined to a width $\Delta x = a_0$. Using the lower bound for the uncertainty principle (and ignoring factors of 2) we find $\Delta p = \hbar / a_0$ (equivalently its wavevector is $k = 1/a_0$). This confinement imparts a kinetic energy $T$ to the electron	
	\begin{align*}
		T \sim \frac{p^2}{m} \sim \frac{\hbar^2 k^2}{m} \sim \frac{\hbar^2}{m a_0^2}
	\end{align*}
	Equating this with one Hartree $e^2/a_0$ yields
	\begin{align*}
	a_0 = \frac{\hbar^2}{me^2}.
	\end{align*}
	
	\item Classically, the nucleus sees the orbiting electron as a current loop with radius $a_0$,
carrying current $I = e/T$ where $T = 2\pi a_0/v$ is the period of the electron's orbit. At
its center, this loop produces magnetic field
	\begin{align*}
	B_N = \frac{2\pi I}{ca_0} = \frac{e}{a_0^2}\frac{v}{c} = \al E_A
	\end{align*}
	
	\item A Bohr magneton is
	\begin{align*}
	\mu_B = \frac{e\hbar}{2mc}  \sim \alpha ea_0.
	\end{align*}
	Equating $\mu_B B_H$ with one Hartree to get
	\begin{align*}
	B_H \sim \frac{1}{\alpha}\frac{e}{a_0^2} = \frac{E_A}{\al}.
	\end{align*}
	
	
	\item See (b) and (c)
	
	\item $B_N$ corresponds to the magnetic field produced by the orbiting electron in a hydrogen atom and as such gives the typical scale of magnetic fields due to motion of atomic
charges. This makes it the preferred choice for an atomic unit of magnetic field.

Since the hydrogen electron moves at velocity $\al$ smaller than the speed of light, the
magnetic field $B_N$ that it produces is a factor of $\al$ times smaller than its electric field
$E_A$. In SI units:
	\begin{align*}
	B_N = \frac{e}{4\pi\epsilon_0} \frac{\al}{c a_0} = 12.52\, \text{T} 
	\end{align*}
	
	At this practically accessible field, the Zeeman interaction with the external field be-
comes comparable to even the strongest magnetic coupling in atoms (L · S coupling).

A Bohr magneton corresponds to the magnetic moment of the hydrogen electron’s
orbital motion. Therefore, $B_H$ corresponds to the magnetic field which interacts with
this electron on the same energy scale as the nuclear electric field. Since the hydrogen
electron moves $\al$ times slower than the speed of light, $B_H \geq 10^5 \text{T}$ is $1/\al$ times larger than $E_A$ and $1/\al^2$ times larger than $B_N$ , i.e. more than 4 orders of magnitude larger than both the atomic and the highest laboratory magnetic fields.
	
	\end{enumerate}



\noindent \textbf{2. Determination of the fine structure constant, $\al$ (Solution by Yu-Kun Lu \& HQB)}


\begin{enumerate}[label=\alph*)]

\item Our goal is to write $\al$ in terms of $f_\infty = cR_\infty$ and $f_e = m_ec^2/h$. This can be done using the definitions. From $\al =e^2/4\pi \epsilon_0 \hbar c$ we find 
	\begin{align*}
	\al^2 = \f{e^4}{(4\pi \epsilon_0)^2 \hbar^2 c^2} = \f{1}{c^2} f_\infty \f{4\pi \hbar}{m_e} = f_\infty \f{2h}{m_ec^2} = \f{2f_\infty}{f_e}
	\end{align*} 
	where we have used the definition for the Rydberg constant $R_\infty$:
	\begin{align*}
	f_\infty = c R_\infty = c\f{m_ee^4}{8 \epsilon_0^2 h^3 c }  =\f{m_ee^4}{8 \epsilon_0^2 h^3 } = \f{m_ee^4}{(4\pi\epsilon_0 \hbar )^2 4\pi \hbar }.
	\end{align*}
	So, 
	\begin{align*}
	\boxed{\al = \sqrt{\f{2f_\infty}{f_e} }}
	\end{align*}
	From here, we see that 
	\begin{align*}
	\al = \sqrt{\f{2cR_\infty}{m_e c^2/h}}
	\end{align*}
	which depends only on the experimental values $R_\infty$ and $h/m_e$. The speed of light $c$ is \textit{defined}. 

\item Typically, this experiment would be done with a non-relativistic (thermal) neutron beam as
its de Broglie wavelength would be much larger and more easily measurable. The momentum
of such a non-relativistic neutron is related by its de Broglie wavelength as $\lambda_B = h/p$. Since $p = mv$, we get $h/m = \lambda_B v$.

\item The photon of frequency $\nu$ carries momentum $p_\gamma = h\nu/c$. In the absorption process the momentum is conserved so
\begin{align*}
mv_R(\nu) = h\nu/c
\end{align*}
or 
\begin{align*}
h/m = cv_R(\nu)\nu.
\end{align*}

\item When the atom absorbs the first photon, it is imparted a recoil velocity $v_R(\nu_1)$ given by the expression from part (c). When the atom emits the photon, it emits it into a beam propagating opposite of its direction of motion. This emission is accompanied by another recoil $v_R(\nu_2)$, in the same direction as the first, corresponding to the frequency $\nu_2$. 

Since the atomic light-recoil velocities are much smaller than the speed of light, their effect
on the atomic resonance is small and $\nu_1 \sim \nu_2 \sim \nu_0$. By the total conservation of energy:
\begin{align*}
0 + h\nu_1 &= h\nu_2 + \frac{m}{2} \lb v_R(\nu_1) + v_R(\nu_2) \rb^2 \\
h(\nu_1 - \nu_2) &= \frac{m}{2}\lp  \frac{h}{cm} \rp^2 (\nu_1 + \nu_2)^2 \\
\frac{\nu_1 - \nu_2}{(\nu_1 + \nu_2)^2} &= \frac{h}{2mc^2}\\
\frac{1}{2} \frac{\Delta \nu c^2}{\nu_0^2} &= \frac{h}{m}.
\end{align*}
Note the factor of $1/2$ which appears due to the second atomic recoil accompanying the
emission of the photon (see PRL \textbf{70}, 18, 2706 (1993)).


\item The relationship between $h/\Delta m$ and $\lambda$ will depend somewhat on the effect of the recoil of the nucleus. If we neglect the recoil, by mass-energy equivalence we get directly
\begin{align*}
\Delta mc^2 = h\nu = \frac{hc}{\lambda} \quad \text{and} \quad  \frac{h}{\Delta m} = \lambda c.
\end{align*}
Since $c$ in SI units is a constant, this only depends on the experimental value of $\lambda$ in
meters. 

Since this is a precision measurement, let’s try to consider the effect of nonrelativistic
nuclear recoil. By mass-energy equivalence and Newtonian momentum conservation:
\begin{align*}
\Delta mc^2 &= \frac{hc}{\lambda} + \frac{(m - \Delta m) v^2}{2}\\
\frac{h}{\lambda} &= (m - \Delta m)v.
\end{align*}
Solving for $\lambda$:
\begin{align*}
\frac{h}{\lambda} &= c\lp \Delta m -m + \sqrt{m^2 - \Delta m^2}\rp \\
\lambda &\approx \frac{h}{c\Delta m} \lb 1 + \frac{\Delta m}{2m } + \mathcal{O}\lp \frac{\Delta m}{m} \rp^2 \rb
\end{align*}

The next-order contribution from the recoil correction scales as $\Delta m/m$. For the 2.2 MeV
$\gamma$-ray transition in deuterium (Nature, \textbf{438}, 1096-1097), the above recoil correction is on the order of $1.2\times 10^{-3}$ -- well inside the accuracy of mass balances. Unfortunately, this correction is off by a factor of two as can be verified by a full relativistic treatment:
\begin{align*}
\frac{hc}{\lambda} + \frac{m'c^2}{1 -  v^2/c^2} &= mc^2 \\
\frac{m'v}{\sqrt{1 - v^2/c^2}} &= \frac{h}{\lambda}.
\end{align*}
Solving for $\lambda$, with $\delta = \Delta m/m$:
\begin{align*}
\lambda &= \frac{h}{2c\Delta m}  \lp 1 + \sqrt{\frac{1 + 3\delta}{1 - \delta}}\rp \\
&\approx \frac{h}{c\Delta m} (1 + \delta + \delta^3  - \delta^4 + 3\delta^5 + \mathcal{O}(\delta^6)  ) \\
&= \frac{h}{c\Delta m} \lp 1 + \frac{\Delta m}{m} + \dots \rp
\end{align*}
This means that the first recoil correction for the emission of a 2.2 MeV $\gamma$-ray from $^2 H$ is on the order of $3 \times 10^{-3}$.

\end{enumerate}

\noindent \textbf{3. Ground state energy of the helium atom (Solution 1 by Yu-Kun Lu)}\\

\begin{enumerate}[label = \alph*)]

\item If we assume the spin wavefunction to be antisymmetric, the spatial electronic wavefunction
will be symmetric. If both electrons are in the $1S$ state, we get
\begin{align*}
\psi(\mathbf{r}_1, \mathbf{r}_2) = \psi_{1S}(\mathbf{r}_1)\psi_{1S}(\mathbf{r}_2).
\end{align*}
The electron-electron interaction energy $\langle e^2/r_{12} \rangle$ can then be written as
\begin{align*}
V_{ee} = \bigg\langle \frac{e^2}{r_{12}} \bigg\rangle  
&= \int d^3 \mathbf{r}_1 \int d^3 \mathbf{r}_2 \, 
\abs{\psi_{1S}(\mathbf{r}_1)}^2 \abs{\psi_{1S}(\mathbf{r}_2)}^2 \frac{e^2}{\abs{\mathbf{r}_1 - \mathbf{r}_2}} \\
&= \int d^3 \mathbf{r}_1 \, \abs{\psi_{1S}(\mathbf{r}_1)}^2 eU(\mathbf{r}_1) \\ 
&= \bra{\psi_{1S}({r}_1)} eU(r_1) \ket{\psi_{1S}(r_2)}
\end{align*}
where
\begin{align*}
U(\mathbf{r}_1) = \int d^3 \mathbf{r}_2 \, \abs{\psi_{1S} (\mathbf{r}_2)}^2 \frac{e}{\abs{\mathbf{r}_1 - \mathbf{r}_2}}
\end{align*}
can be thought of as average electrostatic potential produced by the second electron’s
effective charge distribution $\rho({\mathbf{r}_2})  = e \abs{\psi_{1S} (\mathbf{r}_2)}^2$. If both electrons are in $S$ orbitals, then $\rho(\mathbf{r}_2)$ is spherically symmetric. Then, $U(\mathbf{r}_1)$ can be obtained simply by adding up the contributions due to each spherical shell with charge $dq = e \abs{\psi_{1S}(r_2)}^2 \,4 \pi r_2^2 \,dr_2$:
\begin{align*}
U(r_1) = \int_{r_2 < r_1}  \frac{dq}{r_1}+ \int_{r_2 > r_1} \frac{dq}{r_2}.
\end{align*}
Substituting the hydrogenic ground state wavefunction:
\begin{align*}
\psi_{1S}(r) = \frac{1}{a^{3/2}\sqrt{\pi}} e^{-r/a}, \quad\quad a = a_0/Z
\end{align*}
\begin{align*}
dq &= \frac{e}{\pi a^3}  e^{-2r_2/a} \, 4\pi r_2^2 \, dr_2  \\
U(r_1) 
&= \frac{e}{\pi a^3} 
\lc  
\int_{r_2 < r_1}  \frac{1}{r_1} e^{-2r_2/a} \, 4\pi r_2^2 \, dr_2+ \int_{r_2 > r_1} \frac{1}{r_2} e^{-2r_2/a} \, 4\pi r_2^2 \, dr_2.
\rc\\
&= \frac{e}{r_1} \lb 1 - e^{-2r_1/a }\lp 1 + \frac{r_1}{a}\rp \rb
\end{align*}
Therefore:
\begin{align*}
V_{ee} &= 
\frac{e}{\pi a^3} \int^\infty_0 
\frac{e}{r_1} \lb 1 - e^{-2r_1/a }\lp 1 + \frac{r_1}{a}\rp \rb
e^{-2r_1/a} 4\pi r_1^2 \,dr_1 \\ 
&= \frac{e^2}{a}\int^\infty_0 \frac{1}{x} \lb 1 - e^{-2x}(1+x) \rb e^{-2x} 4x^2 \,dx\\
&= \frac{e^2}{a}\int^\infty_0 \lb e^{-2x} - e^{-4x}(1+x) \rb 4x\,dx \\
&= \frac{5}{8}\frac{e^2}{a} = \frac{5Z}{8} \frac{e^2}{a_0} = 34.014 \text{ eV}.
\end{align*}

Therefore, the first variational estimate for the ground state energy is:
\begin{align*}
E = \lp -Z^2 + \frac{5Z}{8} \rp \frac{e^2}{a_0} = -74.831 \text{ eV}.
\end{align*}


\item In He, the nucleus is  about 4000 times heaver than each electron. Therefore, we neglect nuclear
motion, and rewrite the He Hamiltonian as:
\begin{align*}
H = H_1 + H_2 + \frac{(Z'-2)}{r_1}e^2 + \frac{(Z'-2)}{r_2}e^2 + \frac{e^2}{r_{12}}.
\end{align*}
where
\begin{align*}
H_{1,2} = \frac{p_{1,2}^2}{2m_e} - \frac{Z'e^2}{r_{1,2}}
\end{align*}
correspond to hydrogenic Hamiltonians with reduced mass $m_e$, charge $Z'$ and ground state
eigenfunction $\phi(\mathbf{r})$ with eigenvalue $E_1 = -e^2/2a = -e^2 Z'^2/2a_0$. From class, $\bra{\phi} r \ket{\phi}   = a^{-1} = Z'/a_0$. Finally, from part (a):
\begin{align*}
\bra{\phi(r_1)\phi(r_2)} \frac{e^2}{r_{12}}  \ket{\phi(r_1)\phi(r_2)} = \frac{5Z'}{8} \frac{e^2}{a_0}.
\end{align*}
Therefore,
\begin{align*}
E_0(Z') &= \bra{\phi(r_1)\phi(r_2)} H \ket{\phi(r_1)\phi(r_2)}  \\
&= 2E_1 + \frac{e^2}{a_0}\lb 2(Z'-2) Z' + \frac{5Z'}{8} \rb \\ 
&= \frac{e^2}{a_0} \lp  Z'^2 - \frac{27}{8}Z' \rp.
\end{align*}
$Z'$ is the effective nuclear charge as seen by each electron. The bare nuclear charge is
reduced below $Z = 2$ due to the screening effect of the other electron.

\item  Setting $\p_{Z'} E_0 = 0 \implies Z' = 27/16 =1.6875$. In other words, due to shielding,
each electron sees only 84\% of the nuclear charge. At this optimal $Z'$ value we get the new
best variational estimate of ground state energy:
\begin{align*}
E_0 = -\lp\frac{27}{16}\rp^2 \frac{e^2}{a_0} = -77.489 \text{ eV}.
\end{align*}

\end{enumerate}


\noindent \textbf{3. Ground state energy of the helium atom (Solution 2 by HQB)}


\begin{enumerate}[label=(\alph*)]
	\item Suppose both electrons are in the 1s state and that the spin wavefunction is anti-symmetric, then the two-electron spatial wavefunction is the product of the single-electron wavefunctions $\psi_{100}(\vec{r})$:
	\begin{align*}
	\Psi(\vec{r}_1, \vec{r}_2) 
	&= \psi_{100}(\vec{r}_1)\psi_{100}(\vec{2})\\
	&= \lb \f{1}{\sqrt{\pi}} \lp \f{Z_\text{He}}{a_0} \rp^{3/2}  e^{-Z_\text{He} r_1/a_0}\rb 
	\lb \f{1}{\sqrt{\pi}} \lp \f{Z_\text{He}}{a_0} \rp^{3/2}  e^{-Z_\text{He} r_2/a_0} \rb\\
	&= \f{1}{\pi} \lp \f{Z_\text{He}}{a_0} \rp^3 e^{-Z_\text{He}(r_1 + r_2)/a_0}.
	\end{align*}
	With this, we may calculate the interaction energy $\langle e^2/r_{12}\rangle$ as follows: 
	\begin{align*}
	\bigg\langle \f{e^2}{r_{12}}\bigg\rangle 
	&= \f{1}{\pi^2}\lp \f{2}{a_0} \rp^6 \int d\phi_1 d\theta_1 dr_1    d\phi_2 d\theta_2 dr_2 \, r_1^2\sin\theta_1  r_2^2\sin\theta_2 \, \f{1}{\abs{\vec{r}_1 - \vec{r}_2}}\exp\lb \f{-4 (r_1 + r_2)}{a_0} \rb.
	\end{align*}
	Since we only care about the difference $\vec{r}_1 - \vec{r}_2$, we may take $\vec{r}_1$ to be the $\hat z$-axis for the $\vec{r}_2$ integration. As we integrate through $\vec{r}_1$, this $\hat{z}$-axis changes, but the variables $\theta_1,\phi_1$ simply integrate out. We thus take:
	\begin{align*}
	\abs{\vec{r}_1 - \vec{r}_2} = \sqrt{r_1^2 + r_2^2 - 2r_1r_2\cos(\theta_2)} 
	\end{align*}
	The integral can be rewritten and evaluated in Mathematica (or by hand):
	\begin{align*}
	\bigg\langle \f{e^2}{r_{12}}\bigg\rangle 
	= \f{1}{\pi^2}\lp \f{2}{a_0} \rp^6  \int_{\mathbb{R}^3} d\phi_1 d\theta_1 dr_1 r_1^2\sin\theta_1  \int_{\mathbb{R}^3}   \, 
	\f{d\phi_2 d\theta_2 dr_2 r_2^2\sin\theta_2}{\sqrt{r_1^2 + r_2^2 - 2r_1r_2\cos(\theta_2)}}\exp\lb \f{-4 (r_1 + r_2)}{a_0} \rb
	= \boxed{\f{5e^2}{4a_0}}
	\end{align*}  
	
	What is this energy in eV? Working in SI units, the energy is given by 
	\begin{align*}
	\bigg\langle \f{e^2}{r_{12}}\bigg\rangle = \f{1}{4\pi \epsilon_0} \lp \f{5e^2}{4a_0} \rp \approx 34.014 \text{ eV}. 
	\end{align*}
	The ground state energy is therefore 
	\begin{align*}
	E_g = -108.848 \text{ eV} + 34.014 \text{ eV} = \boxed{-74.831 \text{ eV}}
	\end{align*}
	which is much closer to the measured value of $-79.006 \text{ eV}$.
	
	
	
	
	
	
	
	
	
	Mathematica code:
	\begin{lstlisting}
	(*define \psi_total^2*)
	In[1]:= f = (1/Sqrt[Pi])*(2/a0)^(3/2)*
	Exp[-2 r1/a0]*(1/Sqrt[Pi])*(2/a0)^(3/2)*Exp[-2 r2/a0]
	
	Out[1]= (8 E^(-((2 r1)/a0) - (2 r2)/a0))/(a0^3 \[Pi])
	
	(*Evaluate integral*)
	In[2]:= e^2*
	Integrate[
	f^2*1/Sqrt[r1^2 + r2^2 - 2 r1*r2*Cos[t2]]*r2^2*Sin[t2]*r1^2*
	Sin[t1], {r1, 0, Infinity}, {r2, 0, Infinity}, {t1, 0, Pi}, {t2, 0,
	Pi}, {p1, 0, 2*Pi}, {p2, 0, 2*Pi}]
	
	Out[2]= ConditionalExpression[(5 e^2)/(4 a0), Re[a0] > 0]
	\end{lstlisting}
	
	
	
	
	
	\item Let us now use variational method to obtain the ground state energy for helium, assuming only the interaction energy $\langle e^2/r_{12}\rangle$ is present. To this end, we take the trial wavefunction:
	\begin{align*}
	\psi(\vec{r}_1, \vec{r}_2) = \phi(\vec{r}_1)\phi{\vec{r}_2} = \lb \f{1}{\sqrt{\pi} } \lp \f{Z'}{a_0}\rp^{3/2} e^{-Z'r_1/a_0} \rb \lb \f{1}{\sqrt{\pi} } \lp \f{Z'}{a_0}\rp^{3/2} e^{-Z'r_2/a_0} \rb = \f{1}{\pi}\lp \f{Z'}{a_0} \rp^3 e^{-Z'(r_1 + r_2)/a_0}.
	\end{align*}
	The original Hamiltonian in SI units is 
	\begin{align*}
	\ham = -\f{\hbar^2}{2m}\lp \nabla^2_1 + \nabla_2^2  \rp - \f{e^2}{4\pi \epsilon_0}\lp \f{2}{r_1} + \f{2}{r_2}\rp  +  \f{e^2}{4\pi \epsilon_0}\lp \f{1}{\abs{\vec{r}_1 - \vec{r}_2}}\rp .
	\end{align*}
	Since we are introducing the variational parameter $Z'$, we may rewrite the Hamiltonian to include $Z'$: 
	\begin{align*}
	\ham = \underbrace{-\f{\hbar^2}{2m}\lp \nabla^2_1 + \nabla_2^2  \rp - \f{e^2}{4\pi \epsilon_0}\lp \f{Z'}{r_1} + \f{Z'}{r_2}\rp}_{\ham_0}  +  \f{e^2}{4\pi \epsilon_0}\lp \f{Z'-2}{r_1} + \f{Z'-2}{r_2} + \f{1}{\abs{\vec{r}_1 - \vec{r}_2}}\rp.
	\end{align*}
	which has been intentionally put in the form where the trial wavefunction is an eigenstate of $\ham_0$, which we may refer to as the ``unperturbed'' Hamiltonian. With this, we may evaluate the expectation value of the energy, in Gaussian units\footnote{The conversion to Gaussian units is easy: we simply rescale the Hamiltonian so that energy is counted in units of $e^2/a_0$ -- I won't show the details here.}:
	\begin{align*}
	\langle \ham \rangle 
	&= \langle \ham_0 \rangle + e^2 \bigg\langle \f{Z'-2}{r_1} + \f{Z'-2}{r_2} + \f{1}{\abs{\vec{r}_1 - \vec{r}_2}} \bigg\rangle \\
	&= 2Z'^2 \lp \f{-e^2}{2a_0} \rp + e^2 \lb \f{2Z'(Z'-2)}{a_0} + \f{5Z'}{8a_0}  \rb
	\end{align*} 
	where we have evaluated the terms $\langle 1/r_1\rangle = \langle 1/r_2\rangle$ in Mathematica. The interaction term is also evaluated in Mathematica in a similar way as in the preceeding part. 
	
	
	\textbf{Interpretation:} The free parameter $Z'$ can be thought of as an effective nuclear charge. By having $Z' \neq 2$ we are saying that the electrons influence each other not just through the Coulomb interaction but also through the \textit{shielding} of the nucleus: One electron always sees the nucleus plus a negative cloud due to the other electron, hence it sees an effective nuclear charge of less than 2, as we will see in the next part.
	
	
	Mathematica code:
	\begin{lstlisting}
	(*evaluate integral for <1/r1> = <1/r2> terms*)
	
	In[25]:= e^2*
	Integrate[
	g^2*(Z - 2)/r1*r2^2*Sin[t2]*r1^2*Sin[t1], {r1, 0, Infinity}, {r2, 0,
	Infinity}, {t1, 0, Pi}, {t2, 0, Pi}, {p1, 0, 2*Pi}, {p2, 0, 2*Pi},
	Assumptions -> {Z/a0 > 0}]
	
	Out[25]= (e^2 (-2 + Z) Z)/a0
	\end{lstlisting}
	
	
	\item We find the improved ground state energy variationally:
	\begin{align*}
	E_g = \min_{Z'} \lp -Z'^2 +2Z'(Z'-2) + \f{5 Z'}{8} \rp \f{e^2}{a_0} = -\f{729}{256}\f{e^2}{a_0} \approx \boxed{-77.489 \text{ eV}}
	\end{align*}
	This is attained at $Z' = 27/16 = 1.6875 < 2$.\\
	
	Mathematica code:
	\begin{lstlisting}
	In[27]:= h[Z_] := -Z^2 + (5 Z)/8 + 2 Z (Z - 2);
	
	In[29]:= Solve[D[h[z], z] == 0, z]
	
	Out[29]= {{z -> 27/16}}
	
	In[30]:= h[27/16]
	
	Out[30]= -(729/256)
	
	In[32]:= N[27/16]
	
	Out[32]= 1.6875
	\end{lstlisting}
	
\end{enumerate}





\end{document}








